\chapter{The Path to Collatz}\label{chap:path}

Over the preceding chapters we have accumulated a good deal of machinery: parity
vectors and the $2$ adic conjugacy, martingales\index{martingale} and their failure of uniform
integrability, Lyapunov functions\index{Lyapunov function} and the impossibility theorems for monotone ones, an
optimizer, a cycle bound and a Ces\`aro conjecture\index{Cesaro conjecture@Ces\`aro conjecture}. This chapter steps back and
asks what all of it says about a proof. We set out what a proof of the
conjecture would have to contain, where each piece of the book sits relative to
that, and which statements would close each half.

\section{The conjecture splits in two}\label{sec:ch10:split}

Fix an odd $n\in\N$ and follow its forward orbit under the Syracuse map\index{Syracuse map}. Exactly one of
three things happens. The orbit reaches $1$; or it is unbounded; or it is
eventually periodic on a cycle not containing $1$. The conjecture asserts that
the second and third possibilities never occur, and these are genuinely
separate problems.

The \emph{divergence half} asks whether $\log y_N$ can drift\index{drift} upward forever.
By the Ces\`aro identity\index{Cesaro identity@Ces\`aro identity} (\cref{prop:cesaro}),
\[
  \log y_N-\log y_0=\sum_{j=1}^N\log\Bigl(3+\frac1{y_{j-1}}\Bigr)-N\bar k_N\log2,
\]
so along an unbounded orbit the running mean $\bar k_N$ of the halving
exponents satisfies $\limsup_N\bar k_N\le\log_23\approx1.585$ (\cref{cor:ch9:freq}). Under the
Geometric$(\tfrac12)$ law of \cref{prop:geom} the mean exponent is $2$, which
leaves a heuristic margin of $0.415$; along an unbounded orbit there is no
margin at all, so it describes the model only. The difficulty is that the law describes the
exponent when the argument is drawn uniformly from residue classes, whereas
what is needed is a statement about the exponent when the argument is the
previous iterate of the same map. This is a law of large numbers for a
deterministic sequence.

The \emph{cycle half} asks whether the orbit can close up. A cycle with $r$
odd elements $x_1,\dots,x_r$ and $s$ halvings imposes the exact multiplicative
constraint $\prod_i(3+1/x_i)=2^s$ of \cref{lem:cycleid}: a product of $r$
numbers each slightly above $3$ is exactly a power of $2$. Taking logarithms,
$s\log2-r\log3$ must be positive and smaller than $r/(3M)$, where $M$ is the
least element. No probability is involved; the tools are Diophantine\index{Diophantine approximation},
although Diophantine information alone bounds the length of a cycle without
excluding one (\cref{sec:ch10:cycles}).

The two halves differ in kind, not only in technique. The first is a law of
large numbers we cannot prove. The second is a sequence of finite verifications that we can
carry further and further but never complete by computation alone.

A remark on notation joins the two halves to Part II. The threshold
$\log_23$ that governs divergence equals $1+\beta$, where
$\beta=\log(3/2)/\log2$ is the momentum\index{momentum} coefficient of Switching MM\index{Switching MM}
(\cref{chap:smm}). In the language of the accelerated map\index{accelerated map}, in which one halving
is absorbed into every odd step, the divergence half asks whether the extra
halvings per odd step average more than $\beta$, and the cycle half asks
whether $s/r$ can come within $1/(3M\log 2)$ of $1+\beta$ from above.
By construction the same number sets the inertia of the optimizer; it is also the drift threshold of the orbit and the Diophantine target of the cycle problem.

\section{Closing the divergence half}\label{sec:ch10:divergence}

\Cref{conj:cesaro} is the precise statement required.
It is equivalent to the divergence half, and both directions are immediate (\cref{sec:ch9:conj}): every bounded orbit satisfies it, while by the Ces\`aro identity\index{Cesaro identity@Ces\`aro identity} every unbounded orbit has $\limsup_N\bar k_N\le\log_23$ (\cref{cor:ch9:freq}), so the conjecture holds exactly when no orbit is unbounded.
Its value lies in the mechanism it names.
Three routes suggest themselves, and we show where each fails.

\paragraph{Equidistribution.} Suppose that along an unbounded orbit the iterates were equidistributed over the odd residues modulo $2^m$, for every fixed $m$.
Then $\liminf_N\bar k_N\ge2$ along it, which \cref{cor:ch9:freq} forbids; so if this held along every unbounded orbit there would be no unbounded orbit, and the conjecture would follow.
The hypothesis cannot be asked of every orbit, since an orbit that reaches $1$ is eventually constant.
What is known concerns random starting values rather than a fixed
orbit (\cref{chap:divergence}), which is also how problem O2 of \cref{chap:open} is posed, and the passage from one to the other is the whole
difficulty. For a starting value drawn uniformly modulo $2^m$ the parity vector\index{parity vector} of its first $m$ steps is uniformly distributed over the $2^m$ words (\cref{thm:terras}; Terras\index{Terras, R.} and Everett\index{Everett, C. J.}~\cite{terras1976,everett1977},
see~\cite[Theorem~B]{lagarias1985}); for a random odd starting value its first halving exponents are close in
total variation to independent Geometric$(\tfrac12)$ draws~\cite[Proposition~1.9]{tao2022},
and the deep input in Tao\index{Tao, T.}'s theorem is a fine scale mixing estimate modulo
powers of $3$~\cite[Proposition~1.14]{tao2022}. The finite analogue in
\cref{chap:divergence} makes the point sharply. The model chain on residues
modulo $3^a$, in which the halving exponents are independent draws, mixes\index{mixing} \emph{exactly} in
$a$ strides (\cref{lem:ch9:mixing}), to a law that is far from uniform; but this,
like the rest, is a statement about an ensemble of starting values, averaged by the residue classes, and
none of it constrains the residues visited by one trajectory.

\paragraph{A Lyapunov certificate.} One would like a function $W$ bounded away
from zero and a
constant $\rho<1$ with $W(S(x))\le\rho\,W(x)$ for every odd $x\ge3$, which would
give geometric descent with no probabilistic content. Such a certificate exists
if and only if the conjecture holds (\cref{rem:ch4:exists,rem:ch4:monotone}), so the
question is whether one can be written down without knowing the answer, and
the natural candidates increase with $x$. \Cref{thm:nolyapunov}
shows that no increasing $W$ works: whenever $x\equiv3\pmod4$ the stride ascends, and
half of all odd integers have that form. Composing strides does not help.
\Cref{cor:ch4:multistride} shows that for every fixed horizon $J$ there are
infinitely many $x$ (those congruent to $-1$ modulo $2^{J+1}$) whose first $J$
strides all ascend, so no increasing certificate\index{certificate} over a fixed number of strides exists
either, and a nonincreasing one fails at $x=5$, since $S^{(J)}(5)=1$ for every $J\ge1$.
A certificate must therefore be nonmonotone, or let the horizon depend
on $x$, as the stopping time\index{stopping time} does.

\paragraph{Uniform integrability.} If the martingale $\{x_n\}$ of the
accelerated model, stopped when it first drops to $1$ or below, were uniformly integrable\index{uniform integrability}, optional stopping would give
$\E[x_\tau]=x_0$, which is false since $x_\tau\le1<x_0$. \Cref{prop:uimartingale}
turns this into a proof that uniform integrability fails. In the model the
conditional probability of termination is a uniformly integrable martingale,
but it is the constant $1$ (\cref{prop:absorb}), so there is nothing to
certify. For the deterministic map its counterpart is the indicator of
reaching $1$, which exists trivially; showing that it equals $1$ everywhere is
the conjecture itself (\cref{rem:ch3:cert}).

The first and third routes replace a statement about
a specific trajectory with a statement about an ensemble, and each time the
replacement is the step that cannot be justified; the second avoids the
ensemble, but descent at a fixed horizon fails pointwise for every monotone
$W$. We regard this as the real content of the divergence half. It is not that the heuristic is weak: the model
predicts the total stopping time\index{stopping time} slope to within about two percent and, as
\cref{chap:stochastic} shows, its dispersion to within about ten percent. It is that no mechanism
is known for transferring ensemble statistics to individual orbits in this
setting.

\section{Closing the cycle half}\label{sec:ch10:cycles}

Here the mechanism is better understood, but it can only bound the length of a
cycle, never exclude one. \Cref{lem:cycleid} converts the existence of a cycle
with $r$ odd elements and least element $M$ into the inequality
$M\le r/(3\Lambda_r)$, where $\Lambda_r=(\lceil r\log_23\rceil-r\log_23)\log2$
is $\log2$ times the one sided distance from $r\log_23$ up to the next integer. Two facts then
combine. Exhaustive computation gives a lower bound on $M$, since every integer
below the verified range\index{verification range} reaches $1$ ($2^{68}$ in \cref{thm:cycles}, now
$2^{71}$ by Ba\v{r}ina\index{Barina, D.@Ba\v{r}ina, D.}~\cite{barina2025}). Continued fraction theory gives
an upper bound on $r/(3\Lambda_r)$, since $\Lambda_r$ can be small only when $r$
is a good one sided approximation\index{one sided approximation} denominator of $\log_23$.

The mediant lemma\index{mediant lemma} of \cref{chap:cycles} is what makes the second fact exact.
The one sided distance $\delta(r)=\lceil r\lt\rceil-r\lt$ is additive
modulo $1$, so if a smaller $r'$ has $\delta(r')<\delta(r)$ then $\delta(r)$
splits as $\delta(r')+\delta(r-r')$ and $\delta(r)/r$ is a mediant of two
smaller ratios.
It follows that the first $r$ not excluded by the verified range, and every record of $r/\delta(r)$, is a record of $\delta$, and \cref{thm:ch8:onesided} identifies these records with the rungs of the semiconvergent\index{convergent}\index{semiconvergent} ladder, which is the classical description of the best one sided approximations
(Han\v{c}l\index{Hancl, J.@Han\v{c}l, J.} and Turek\index{Turek, O.}~\cite[Theorem~4.5(ii)]{hancl2019}). There are only $53$ rungs below $2\times10^{13}$, the bound is
evaluated at each, and \cref{thm:cycles} follows without any unverified
proposition: every nontrivial cycle has $r\ge72\,057\,431\,991$ odd elements.

The gap is that this excludes finitely many $r$ at a time. A large enough improvement in
the verification range moves the exclusion forward to a later rung, and the
excluded range grows roughly like the square root of the verified range:
$2^{40}$ gives $r\ge1.08\times10^{7}$ and $2^{68}$ gives $r\ge7.2\times10^{10}$. The table in \cref{chap:cycles} shows the steps: verification
to $2^{71}$~\cite{barina2025} gains nothing over $2^{68}$ for this inequality, because the next obstruction sits
at $2^{71.88}$; verification to $2^{72}$ would move the bound to
$137\,528\,045\,312$, the current record, which Hercher\index{Hercher, C.}'s sharper accounting
of the cycle already obtains from $3\cdot2^{69}<2^{71}$~\cite[Corollary~29]{hercher2023};
and then nothing further happens until beyond $2^{75}$.
Can a bound on $\Lambda_r$ valid for all $r$ at once, an effective
irrationality measure\index{irrationality measure} for $\log_23$, close the gap? It cannot.
\Cref{rem:irrmeasure} shows that a measure $\nu$ with constant $c$ yields only
$M\le r^{\nu}/(3c\log2)$, which grows at least like $r^2$, and
\cref{prop:ch8:insufficient} shows more: for every verification range there
are infinitely many $r$ that the Diophantine\index{Diophantine approximation} inequality cannot exclude, whatever
is known about $\log_23$. Even an idealized measure with $\nu=2$ and the best
constant consistent with the computed partial quotients is beaten by the exact
ladder computation by a factor of twenty at $2^{68}$.

So the cycle half is \emph{not} a question about one real number. What is missing is
arithmetic information about the parity word. The window lemma\index{window lemma} of
\cref{chap:cycles} makes this precise. The affine identity of
\cref{chap:parity} writes every element of a cycle as $c_k(x)/(2^k-3^m)$, with
$c_k(x)$ a sum of terms $3^j2^i$ indexed by the parity word, and the trivial
bound $c_k\ge3^m-2^m$ gives a lower bound for the least element which, together
with the upper bound $M<r/(3\Lambda)$, confines $M$ to a window whose ends are
a factor of about $r/3$ apart.
An arithmetic improvement of the upper end by a factor slightly more than $r/3$, together with a verification up to a fixed bound, would exclude every nontrivial cycle (\cref{cor:ch8:window}); raising the lower end alone empties the window only for cycles with small $\Lambda$.
The exploration of parity words in
\cref{chap:cycles} suggests why this is hard: for balanced words\index{balanced words} the least element
sits, in our computations, at a nearly fixed fraction, about $0.72$, of the top of the window, and the
residues of $c_k$ modulo $2^k-3^m$ look uniformly random. A proof must exploit
divisibility. The Erd\H{o}s $125$ analogy suggests where to look: there the good
approximations $3^m\approx4^k$ thin the sumset and drive its lower density to
zero; for cycles, approximations of the same constant, $2^k$ slightly above $3^m$, are what a cycle needs, and the
question is whether the digit structure of the word prevents
$2^k-3^m$ from ever dividing $c_k$.

\section{Where the optimization framework fits}\label{sec:ch10:fits}

The framework does not prove either half. What it supplies is the vocabulary in
which the halves separate cleanly, and it collects in one place the constants a
proposed mechanism must reproduce: the drift\index{drift} $\ds=\log(\sqrt3/2)$, the step variance
$\sigma^2\approx0.302$, the total stopping time\index{stopping time} slope $1/|\ds|\approx6.95$, the stride
constants $\log(5/3)$ and $|\log(3/4)|$ of \cref{def:nesterov}, the parity autocorrelations\index{autocorrelation} $(-\tfrac12)^h$ and the long run variance factor $\tfrac13$ of the model for the standard map $T_0$ (\cref{thm:autocorr,cor:longrun}). These matter for a
practical reason. When one proposes a mechanism for the conjecture, the first
test is whether it reproduces the known constants. A mechanism that predicts a
drift other than $\log(\sqrt3/2)$ per step, or a total stopping time slope other than about
$6.95$, for the accelerated map\index{accelerated map}, can be discarded without further work.

\begin{sloppypar}
It also records negative results. \Cref{thm:nolyapunov} and
\cref{cor:ch4:multistride} rule out monotone certificates\index{certificate} over a bounded horizon,
\cref{prop:mmfails} rules out the natural majorization\index{majorization} scheme, and
\cref{prop:uimartingale} rules out the naive martingale\index{martingale} argument. The pitfall remarks scattered through the book
rule out a further collection of arguments that look like progress. The
Collatz iteration is not a descent method, and no monotone reparametrization
makes it one.
\end{sloppypar}

And the framework gave something back. The arithmetic of the map, read as an
optimizer, produced Switching MM\index{Switching MM}, a method whose structure (cyclic step\index{cyclic step sizes}s of
sizes $\alpha$ and $(1+\beta)\alpha$ on quadratics, \cref{thm:smmquad}) and
whose noise behaviour (amplification bounded by $5/3$ for every $\beta<1$,
\cref{thm:smmvar,prop:ch7:bound,prop:ch7:uniform}) can be analysed exactly. The constant $\beta$ does not prove
anything about Collatz orbits; it is simply a well motivated, tuning free
default, and \cref{chap:noise} is careful about how much that is worth.

\section{Verification as a partner}\label{sec:ch10:verification}

A last theme runs through Part III. The cycle bound is assembled from three
steps: the affine identity, the cycle inequality $3^m<2^k$ with its upper
companion, and a finite computation with continued fractions\index{continued fraction}. The first step
and the lower half $3^m<2^k$ of the second are now machine checked in Lean\index{Lean}
(\cref{chap:cycles}); the upper half is an inequality between natural numbers
within reach of the same tools; and the mediant lemma\index{mediant lemma} and \cref{thm:ch8:onesided} reduce the third to a
short certificate\index{certificate}: a list of ladder rungs $u$ and numerators $\lceil u\lt\rceil$ together with
rational bounds on $\log_23$ that verify each inequality by interval
arithmetic. Formalizing the upper half, the mediant lemma, \cref{thm:ch8:onesided} and a verified squaring routine, and replaying that certificate in Lean, would make the cycle bound machine checked, conditional only on the verification below $2^{68}$ (project P1 of \cref{chap:open}). The
divergence half offers no such target, and the contrast is itself informative:
every finite part of the cycle half is within reach of verification, while the
divergence half, which requires a law of large numbers for one orbit, is not.

The next chapter takes up the computational side of this accounting: which
parts of the programme computation can finish, which it provably cannot, and
what the optimization reading says about where to look.

\section*{Notes}

The split into a divergence half and a cycle half is classical; see the surveys
of Lagarias\index{Lagarias, J. C.}~\cite{lagarias1985,lagarias2010overview}. The observation that the probabilistic attacks
on the divergence half replace an orbit by an ensemble is implicit in the
parity heuristic of Crandall\index{Crandall, R. E.}~\cite{crandall1978} and the stochastic models of Lagarias and Weiss\index{Weiss, A.}~\cite{lagarias1992} and explicit in the
discussion surrounding Tao\index{Tao, T.}'s theorem~\cite{tao2022}. The reduction of the cycle
half to Diophantine\index{Diophantine approximation} approximation of $\log_23$ goes back to the work of
Crandall~\cite{crandall1978}, Steiner\index{Steiner, R. P.}~\cite{steiner1977} and Eliahou\index{Eliahou, S.}~\cite{eliahou1993}. That it
cannot suffice is a point Halbeisen\index{Halbeisen, L.} and Hungerb\"uhler\index{Hungerbuhler, N.@Hungerb\"uhler, N.}~\cite[\S5]{halbeisen1997}
made in another form, since rational cycles satisfy the same inequalities; the
proof of \cref{prop:ch8:insufficient}, and the window corollary, are ours.

\section*{Exercises}

\begin{exercise}
Show that the three possibilities for a Syracuse orbit (reaching $1$, being
unbounded, being eventually periodic on a cycle not containing $1$) are
mutually exclusive and exhaustive.
\end{exercise}

\begin{exercise}
Show that along a cycle with $r$ odd elements and $s$ halvings the averages
$\bar k_N$ converge to $s/r$, and that $0<s/r-\log_23<1/(3M\log2)$. Conclude that
the Ces\`aro condition is satisfied by every cycle, so that no argument based on
Ces\`aro averages alone can exclude cycles.
\end{exercise}

\begin{exercise}
Suppose an effective irrationality measure\index{irrationality measure} $\nu=3$ with constant $c=10^{-3}$
were known for $\log_23$. Combine \cref{rem:irrmeasure} with the verified range\index{verification range}
$X=2^{68}$ to show that it excludes exactly the cycle lengths
$r<(3cX\log2)^{1/3}$, about $8.5\times10^{5}$, far fewer than
\cref{thm:cycles}. Explain why no choice of $(\nu,c)$ can exclude all $r$.
\end{exercise}

\begin{exercise}
Show that if $\liminf_N y_N<\infty$ along a Syracuse orbit then the orbit is
eventually periodic. Deduce that an orbit is unbounded if and only if
$y_N\to\infty$, so that the divergence half concerns orbits that escape to
infinity, not merely orbits that are unbounded along a subsequence.
\end{exercise}

\begin{exercise}
List the constants of \cref{sec:ch10:fits} and, for each, state one proposed
mechanism for the conjecture that it would falsify if the mechanism predicted a
different value.
\end{exercise}
